\section{Related work}

\subsection{Censored environmental data}
Concentrations reported below a reporting limit are \emph{left-censored}.
Substituting a constant for them biases means, variances and trends by an
amount that depends on the censoring rate
\cite{helsel2006fabricating,helsel2012nondetects}, and Antweiler and Taylor
show it directly by treating coincident uncensored data as ground
truth~\cite{antweiler2008censored}. Robust alternatives -- regression on order
statistics, maximum likelihood estimation, Kaplan--Meier -- exist and go unused
in routine practice~\cite{shoari2018nondetects}. Our concern is upstream of all
of them: censoring status is discarded before the analyst opens the table, and
no estimator recovers what the schema could not hold.

\subsection{Observation vocabularies}
In SOSA/SSN \cite{janowicz2019sosa,haller2019ssn}, \texttt{sosa:Observation}
carries \texttt{hasSimpleResult}, a point value;
\texttt{ssn-system:DetectionLimit} is a property of a \emph{sensor}, never of an
observation or its result, so nothing can state that a measurement fell below
it. SOSA's \texttt{Sampling}, \texttt{Sampler} and \texttt{Sample} do not link
an analysis to a detection limit either, so the absence is not repaired by
choosing the sampling half of the vocabulary over the sensing half.
Section~\ref{sec:results-gap} shows why: the same absence appears in an
analytical chemistry ontology that has no sensor concept in it at all. What
these vocabularies share is not a model of measurement but a placement --- the
limit is held on whatever produced the result, and never on the result.

\paragraph{Why extend SOSA rather than something else.}
The choice of base is itself a claim, so we state the reasoning. SOSA/SSN is a
joint W3C Recommendation and OGC standard, which no alternative in this space
matches, and it is what the monitoring vocabularies compared in
Section~\ref{sec:results-gap} are built on --- so extending it is what makes the
diagnosis actionable for them rather than an argument for a parallel stack. We
add classes and properties and redefine no SOSA term. Three alternatives were
considered. \emph{O\&M / OMS} (ISO~19156), of which SOSA is the lightweight
\textsc{owl} rendering, does carry a \texttt{resultQuality} slot; but it is
typed to accept any data-quality element, so it is somewhere to put a string
rather than a modality a query can return, and the distinction between a bound
and a measurement is exactly what an unconstrained slot fails to preserve.
\emph{QUDT} and \emph{OM} model units and quantities, which CENSO reuses rather
than replaces. A \textsc{bfo}-grounded design in the \textsc{obo} style would
buy rigour at a cost that lands precisely on this subject matter: ``the record
cannot decide'' is an epistemic state of a document, not an entity in the world,
and representing absence and non-findings is a long-standing difficulty in that
tradition. The commitment we make instead is narrower and checkable --- the
vocabulary stays inside OWL~2~RL --- verified against the profile grammar
rather than assumed, which Section~\ref{sec:results-profile} reports --- and
what the profile cannot express is stated
and moved to SHACL rather than asserted and left unenforced.

\subsection{Water-quality exchange standards}
\label{sec:related-exchange}
Censoring is not unrepresented in water-quality data exchange, and the claim
this paper makes has to be stated against the standards that do carry it. OGC
WaterML~2.0~\cite{ogc2014waterml2} attaches a \texttt{censoredReason} to a
time--value pair and admits a censoring threshold as an untyped point-level
\texttt{qualifier}; ODM2~\cite{horsburgh2016odm2} and CUAHSI WaterML~1.1 carry a
censor code (less-than, non-detect, present but not quantified) on every value;
the SeaDataNet measurand qualifier flags~\cite{nerc_l20}, published as SKOS,
separate a value below detection from one below the limit of quantification;
and the US EPA Water Quality Exchange schema~\cite{epa_wqx3} goes furthest,
recording a detection condition and, on each result, any number of typed
detection and quantitation limits. \textbf{A censored result with its own limit
is therefore expressible today}, in XML and relational exchange formats
(Table~\ref{tab:exchange}). Three
things are not. In the SeaDataNet flags and in ODM2 the limit occupies the value
slot, so a reported number and the limit it was measured against cannot both be
kept. None of these standards is an ontology in which the limit, the result and
the applicable standard are related by axioms a reasoner or a shape can act on:
WQX can list a water-quality criterion among a result's limits, but as one more
limit type and not as a comparison. And none has a compliance outcome, so the
third outcome Article~3(3b) requires, and the reason it was reached, has nowhere
to go. The claim of this paper is confined accordingly. It is not that a
non-detection cannot be flagged; it is that no surveyed vocabulary or exchange
standard provides, in RDF/OWL, a result-level bound tied to the limit it came
from and to a three-valued verdict against a conditional standard.

\subsection{Laboratory and analytical vocabularies}
The concepts CENSO adds would most plausibly already exist not in a water
ontology but in a laboratory one, so the comparison in
Section~\ref{sec:results-gap} includes them. \textbf{CHMO, the \textsc{obo}
Chemical Methods Ontology, does declare \texttt{limit of detection} and
\texttt{limit of quantification}, both defined from the IUPAC Gold Book.} It
binds them to the method, as a \texttt{figure of merit} of an assay, and has no
property relating a limit to a result --- so, as with the sensor case, the limit
is known about the procedure and forgotten about the measurement. AFO, built by
a pharmaceutical consortium for analytical instrument data, contains no
detection or quantification limit term at all; its \texttt{lower bound} and
\texttt{upper bound} are \textsc{bfo} roles rather than values on a result. In
STATO, where censoring is a routine concept in survival analysis, it appears
only in prose. The gap is therefore not an artefact of having looked only at
water: the vocabulary that defines the limit most carefully still cannot say
that a given result fell below one.

\subsection{Water domain ontologies}
Water-quality vocabularies are plentiful~\cite{water_onto_review}: WaWO+ for
river basins~\cite{olivafelipe2017wawoplus}, OPO~\cite{wang2020opo},
InWaterSense~\cite{jajaga2015expert,ahmedi2013wqm},
SAREF4WATR~\cite{etsi2025saref4watr}, the Water Health Open Knowledge Graph
(WHOW)~\cite{whow2023onto,whokg2025}, query technologies over
SSN~\cite{calbimonte2012ssw} and water-treatment
interoperability~\cite{saka2026acquirium}.

Treating a regulation as machine-readable rules rather than as numbers in a
column is not new outside this domain: compliance checking over building codes
has been built on the same idea~\cite{bouzidi2012compliance}. What that
literature does not face is a measurement whose result is a bound, so the
verdict it computes is always determinate.

WHOW is the closest prior work and, unlike most, publishes at a dereferenceable
IRI. Its \texttt{Range} class bounds an \texttt{ObservationValue}, reading
\texttt{<10} as $[0,10]$. \textbf{WHOW therefore does represent a censored
result as an interval}: this work extends it rather than filling a void. What it
does not record is why the interval exists: it does not separate censoring from
spread, does not tie the interval to its limit, and cannot pose an indeterminate
compliance outcome. Section~\ref{sec:results-gap} compares \num{21} ontologies;
none carries the epistemic status of the measurement, and
Section~\ref{sec:results-separation} shows what follows --- on one real record,
\num{20} of the \num{21} cannot tell a non-detection from a measurement
reporting the same number.

\paragraph{Analytical coverage already changes conclusions.}
Moschet et al.~\cite{moschet2014screening} find that the pesticides a routine
programme omits dominate the risk picture; Malaj et
al.~\cite{malaj2014continental} read continental-scale risk to freshwater
ecosystems from the Waterbase records of Section~\ref{sec:results-waterbase};
Kase et al.~\cite{kase2018estrogens} report estrogen standards below what
routine laboratories can quantify at all. They ask what better chemistry would
reveal; we ask what the record establishes as it stands. The answer is often
nothing, and no substitution rule repairs it. Ahkola et
al.~\cite{ahkola2024uncertainty} treat uncertainty as a quantity to estimate and
propagate; where the threshold falls inside $[0,\mathrm{LOD}]$ the measurement
did not discriminate, so no distribution can be placed over the verdict -- a
representational gap, not an estimation problem.

\paragraph{A fourth reason the Directive misses risk.}
Weisner et al.~\cite{weisner2022threereasons} give three reasons the Water
Framework Directive fails to identify pesticide risk: an outdated analyte list,
biased sampling sites, and thresholds outside the Directive going unconsidered.
Each concerns \emph{which} substances are looked for and \emph{where}; ours
applies to the listed substances at monitored sites, where the record often
cannot decide the question it was collected to answer. Coquery et
al.~\cite{coquery2005analytical} raised uncertain analytical capability when the
priority substance list was first implemented; the Commission's Watch List
review restates it~\cite{loos2024watchlist}. What is missing is a way to
\emph{count} the consequence, which requires representing the failure rather
than substituting past it.
